% ch09-analytics-summary
\section{التحليلات والخلاصة وأسئلة المناقشة}

وصلنا للفصل الأخير. في الفصول اللي فاتت مشينا خطوة بخطوة: بنينا الطبقات
(\en{Clean Architecture})، وفهمنا إزاي كل طلب بيتحوّل لأمر أو استعلام
(\en{CQRS})، وأمّنّا الدخول، وركّبنا مزايا الـ\en{AI}، وشبكنا الفيديو والدفع
والمهام الخلفية. في الفصل ده هنعمل حاجتين: الأولى نشرح \textbf{خطّ التحليلات}
(إزاي بيانات النظام الحيّة بتتحوّل لتقارير ولوحات تحكّم)، والثانية \textbf{نلمّ
الكتاب كلّه} في مكان واحد، وبعدها نجهّز \textbf{أسئلة لجنة المناقشة} بإجاباتها
عشان تبقى جاهز في يوم الدفاع.

هدفنا هنا مش نضيف تعقيد جديد، لكن نربط الصورة. فخُد نفَس، وتعالَ نمشي واحدة واحدة.

% ═══════════════════════════════════════════════════════════════════════════
\subsection{خطّ التحليلات: من قاعدة البيانات الحيّة إلى لوحات \en{Power BI}}
% ═══════════════════════════════════════════════════════════════════════════

النظام بيولّد بيانات كتير كل يوم: طلبات منح، حجوزات استشاريين، مدفوعات، تقييمات.
قاعدة البيانات التشغيلية (\en{SQL Server}) شغلها الأساسي إنها تخدم المستخدمين بسرعة،
مش إنها تعمل تقارير تحليلية ثقيلة. عشان كده بنفصل عالَم \textbf{التشغيل} عن عالَم
\textbf{التحليل}، وده اللي بنسمّيه \textbf{خطّ التحليلات} (\en{Analytics Data
Pipeline})، وهو مسؤولية الوحدات \en{PB-015} لحد \en{PB-018}.

\begin{idea}[الفكرة: يعني إيه \en{Data Warehouse}؟]
مخزن البيانات (\en{data warehouse}) هو قاعدة بيانات \textbf{منفصلة} مبنيّة مخصوص
للتحليل والتقارير، مش للمعاملات اللحظية. بننقل نسخة من بيانات النظام إليه، ونعيد
ترتيبها بشكل يخلّي الأسئلة التحليلية (مثال: «كام طلب منحة اتقدّم الشهر ده لكل
دولة؟») تتجاوب بسرعة، من غير ما نضغط على قاعدة البيانات الحيّة اللي بيشتغل عليها
الطلاب في نفس اللحظة.
\end{idea}

\subsubsection{المراحل بالترتيب}

الخطّ عندنا بيمشي في خمس محطات أساسية، من الشمال لليمين:

\begin{idea}[الفكرة: يعني إيه \en{CDC}؟]
\en{Change Data Capture} (اختصارًا \en{CDC}) معناه «التقاط التغييرات فقط». بدل ما
ننسخ الجدول كامل كل ليلة (تبذير)، الـ\en{CDC} بيراقب قاعدة البيانات ويلتقط
\textbf{الصفوف اللي اتغيّرت بس} — أُضيفت أو اتعدّلت أو اتمسحت — وينقلها. ده بيخلّي
النقل خفيف وسريع.
\end{idea}

\begin{idea}[الفكرة: يعني إيه \en{Star Schema}؟]
المخطّط النجمي (\en{star schema}) هو طريقة لترتيب جداول التحليل. في النص جدول
واحد كبير اسمه \textbf{جدول الحقائق} (\en{fact table}) بيحتوي الأحداث القابلة
للقياس (مثال: كل عملية دفع بمبلغها وتاريخها)، وحواليه \textbf{جداول الأبعاد}
(\en{dimension tables}) اللي بتوصف السياق (الطالب، الدولة، نوع المنحة، الزمن).
الشكل بيطلع زي النجمة: حقيقة في النص وأبعاد على الأطراف. الترتيب ده بيخلّي التقارير
سهلة وسريعة جدًّا.
\end{idea}

نبدأ من \en{SQL Server} (المصدر التشغيلي). الـ\en{Azure Data Factory} بيسحب
التغييرات بتقنية \en{CDC} وينزّلها في طبقة اسمها \textbf{\en{Bronze}} — دي نسخة
خام طبق الأصل من البيانات كما هي. بعدها أداة \en{dbt} بتنضّف وتحوّل: من
\en{Bronze} لطبقة \textbf{\en{Silver}} (بيانات متنظّفة وموحّدة)، ومنها لطبقة
\textbf{\en{Gold}} (الـ\en{star schema} النهائي الجاهز للتحليل). بعد كده
\en{Synapse Serverless} بيعرض الـ\en{Gold} على هيئة \en{views} يقدر أي حد
يستعلم عنها بلغة \en{SQL} من غير سيرفر دائم شغّال. وأخيرًا \en{Power BI} بيرسم
اللوحات التفاعلية اللي المشرف بيبصّ عليها.

\begin{center}
\begin{tabular}{L{3.1cm} L{4.0cm} L{5.2cm}}
\toprule
\textbf{\color{spPrimary}المرحلة} & \textbf{\color{spPrimary}الأداة} & \textbf{\color{spPrimary}بتعمل إيه} \\
\midrule
المصدر التشغيلي & \en{SQL Server 2022} & قاعدة البيانات الحيّة اللي فيها كل المعاملات \\
الالتقاط والنقل & \en{Azure Data Factory} (\en{CDC}) & تلتقط الصفوف المتغيّرة فقط وتنقلها \\
طبقة \en{Bronze} & \en{Data Lake} (خام) & نسخة طبق الأصل من البيانات كما هي \\
التحويل & \en{dbt} (\en{Silver} ثم \en{Gold}) & تنضّف وتوحّد ثم تبني الـ\en{star schema} \\
طبقة الاستعلام & \en{Synapse Serverless} & تعرض الـ\en{Gold} كـ\en{views} بلا خادم دائم \\
العرض & \en{Power BI} & لوحات تحكّم تفاعلية للمشرفين \\
الزمن الحقيقي & \en{Event Hubs} ثم \en{Stream Analytics} & يبثّ الأحداث لحظيًّا ويكشف الشذوذ \\
التنبيه & \en{PagerDuty} & ينبّه الفريق فورًا عند اكتشاف شذوذ \\
\bottomrule
\end{tabular}
\end{center}

\subsubsection{الفرع اللحظي: كشف الشذوذ في الوقت الحقيقي}

الخطّ اللي فوق ده \textbf{دفعيّ} (\en{batch}): بيشتغل مرّة كل ليلة. بس فيه أحداث
لازم ننتبه لها \textbf{في اللحظة}، مش بكرة الصبح — زي طفرة مدفوعات مفاجئة ممكن
تكون احتيال. عشان كده فيه \textbf{فرع لحظي} موازي: النظام بيبثّ ثلاثة أحداث
عالية القيمة إلى \en{Azure Event Hubs} أول ما تحصل، و\en{Azure Stream Analytics}
بيقرأها ويكشف الشذوذ، ولو لقى حاجة غريبة بيبعت تنبيه لـ\en{PagerDuty} وكمان
بيغذّي بلاطة حيّة في \en{Power BI}.

\begin{keyp}[نقطة مهمة: التحليلات لا تعطّل التشغيل أبدًا]
مبدأ التصميم هنا إن خطّ التحليلات \textbf{انتهازيّ} (\en{opportunistic}): لو
وقع، النظام الأساسي بيفضل شغّال عادي. الطالب اللي بيدفع أو بيقدّم طلب مش هيتأثّر
لو الـ\en{Event Hubs} كان واقع. الحدث الضايع بيتعوّض بالحدث اللي بعده، مش بإيقاف
المستخدم.
\end{keyp}

تعالَ نشوف ده في الكود الحقيقي. الأحداث الثلاثة اللي بنبثّها هي:
\en{ApplicationSubmitted} (تقديم طلب)، و\en{PaymentCaptured} (تحصيل دفعة)،
و\en{BookingCompleted} (اكتمال جلسة). كل واحد ليه مُعالِج
(\en{MediatR notification handler}) رفيع بيبني ظرفًا صغيرًا ويسلّمه للناشر.

\textbf{هنشوف إيه:} مُعالِج حدث «تحصيل دفعة» وهو بيوجّه الحدث لخطّ التحليلات
(من ملف \code{DomainEventPublishingHandlers.cs}).

\begin{english}
\begin{lstlisting}[language={[Sharp]C}]
public sealed class PaymentCapturedStreamHandler(
    IEventPublisher publisher,
    ILogger<PaymentCapturedStreamHandler> logger)
    : INotificationHandler<PaymentCapturedEvent>
{
    public async Task Handle(PaymentCapturedEvent n, CancellationToken ct)
    {
        // Forward the domain event to the streaming pipeline
        await publisher.PublishAsync(
            "PaymentCaptured",
            new
            {
                paymentId   = n.PaymentId,
                type        = n.Type.ToString(),
                amountCents = n.AmountCents,
                payerUserId = n.PayerUserId,
                payeeUserId = n.PayeeUserId,
            },
            ct).ConfigureAwait(false);
    }
}
\end{lstlisting}
\end{english}

نلاحظ حاجات بسيطة: الصنف بيطبّق \code{INotificationHandler<PaymentCapturedEvent>}،
يعني \en{MediatR} بينده عليه أوتوماتيك أول ما الحدث يتنشر بعد نجاح المعاملة. جوّه
\code{Handle} مفيش أي منطق أعمال — بس نداء واحد لـ\code{publisher.PublishAsync}
بنبعت له اسم الحدث \code{"PaymentCaptured"} وكائن صغير فيه الأرقام المهمّة
(المبلغ بالسنت، ونوع الدفعة، والدافع، والمستفيد). المُعالِج مقصود إنه رفيع كده،
وكل التعامل مع الأخطاء متروك للناشر.

\textbf{هنشوف إيه:} الناشر الحقيقي (\en{Azure Event Hubs}) وإزاي بيبلع الأخطاء
عشان ميعطّلش المعاملة (مقتطف مختصر من \code{EventHubPublisher.cs}).

\begin{english}
\begin{lstlisting}[language={[Sharp]C}]
public async Task PublishAsync(string eventType, object payload, CancellationToken ct = default)
{
    try
    {
        var envelope = new
        {
            eventType,
            occurredAt = DateTimeOffset.UtcNow,
            source = "ScholarPath.API",
            data = payload,
        };
        var json = JsonSerializer.Serialize(envelope);
        var data = new EventData(Encoding.UTF8.GetBytes(json));

        await using var producer = new EventHubProducerClient(
            _opts.ConnectionString, _opts.HubName);
        using var batch = await producer.CreateBatchAsync(ct);
        batch.TryAdd(data);
        await producer.SendAsync(batch, ct);        // send to Azure Event Hubs
    }
    catch (Exception ex)
    {
        // A hub outage must NOT roll back the business transaction
        logger.LogError(ex, "EventHub: failed to publish {EventType}; continuing.", eventType);
    }
}
\end{lstlisting}
\end{english}

الفكرة المحوريّة هنا سطرين: أول حاجة بنغلّف الحمولة في \textbf{ظرف} (\en{envelope})
فيه نوع الحدث ووقت حدوثه ومصدره، ونحوّله \en{JSON} ونبعته لِـ\en{Event Hubs}.
تاني حاجة — وهي الأهم — الـ\code{catch} بيمسك أي استثناء \textbf{ويسجّله بس من
غير ما يعيد رميه}. ده اللي بيحقّق مبدأ «التحليلات ما تعطّلش التشغيل»: لو الـ\en{hub}
واقع، الدالة بتكمل بهدوء والمعاملة الأصلية بتنجح عادي.

\textbf{هنشوف إيه:} إزاي النظام بيختار بين الناشر الحقيقي والبديل الآمن حسب
الإعداد (مقتطف من \code{DependencyInjection.cs}).

\begin{english}
\begin{lstlisting}[language={[Sharp]C}]
// Real publisher only when a connection string is configured; otherwise a stub
var ehConnString = config.GetValue<string>($"{EventHubOptions.SectionName}:ConnectionString");
if (!string.IsNullOrWhiteSpace(ehConnString))
    services.AddSingleton<IEventPublisher, EventHubPublisher>();
else
    services.AddSingleton<IEventPublisher, StubEventPublisher>();
\end{lstlisting}
\end{english}

ده نمط بنكرّره في كل النظام: لو الإعداد موجود (سلسلة اتصال حقيقية) نركّب الخدمة
الحقيقية \code{EventHubPublisher}؛ ولو مش موجود (زي بيئة التطوير على جهاز مطوّر)
نركّب بديلًا آمنًا \code{StubEventPublisher} بيسجّل الأحداث ويكمّل من غير ما يرمي
خطأ. النتيجة: المطوّر يقدر يشغّل النظام كامل من غير أي حساب \en{Azure}.

\begin{secnote}
كل الأحداث المبثوثة \textbf{ما فيهاش} بيانات حسّاسة زايدة: بنبعت مُعرِّفات
(\en{IDs}) وأرقام مجمّعة زي المبلغ والنوع، مش أسماء ولا إيميلات. كده لو اتّسرب
الـ\en{stream} ما يكشفش هويّات المستخدمين — وده يتماشى مع نهج تقليل البيانات
(\en{data minimisation}) اللي شفناه في فصل الأمان.
\end{secnote}

% ═══════════════════════════════════════════════════════════════════════════
\subsection{ملخّص الكتاب كلّه: ماذا تعلّمنا فصلًا فصلًا}
% ═══════════════════════════════════════════════════════════════════════════

خلّينا نلمّ الرحلة. كل فقرة تحت بتلخّص فصلًا كامل في سطرين:

\textbf{الفصل التمهيدي —} عرّفنا \en{ScholarPath}: منصّة منح دراسية ثنائية اللغة
(عربي وإنجليزي بدعم \en{RTL} كامل)، بتخدم الطلاب والاستشاريين والجهات المانحة
والمشرفين، ومقسّمة إلى \textbf{ثماني عشرة وحدة} (\en{PB-001..PB-018}) بتغطّي دورة
الحياة كاملة من إنشاء الحساب لحد تخزين البيانات للتحليل.

\textbf{فصل \en{Clean Architecture} —} اتعلّمنا إن الكود مقسّم لطبقات متداخلة،
والاعتماديّات بتشاور \textbf{للداخل بس}: طبقة \en{Domain} (قلب النظام) ما تعرفش
حاجة عن قاعدة البيانات ولا الإطار ولا الواجهة. ده بيخلّي منطق الأعمال قابل
للاختبار والنقل، وبيسمح باستبدال أي خدمة خارجية ببديل من غير ما نلمس القلب.

\textbf{فصل \en{CQRS} والسلامة —} كل عمليّة بقت إمّا \textbf{أمرًا} (\en{Command})
بيغيّر الحالة، أو \textbf{استعلامًا} (\en{Query}) بيقرأ بس. \en{MediatR} بيوصّل
كل طلب لِمُعالِجه، والمتحكّمات (\en{Controllers}) بقت رفيعة جدًّا. وسلوكيّات
الأنبوب (\en{pipeline behaviours}) زي التحقّق والتدقيق بتلفّ كل مُعالِج تلقائيًّا.

\textbf{فصل المصادقة —} الدخول بيعتمد \en{ASP.NET Core Identity} مع رموز
\en{JWT} موقّعة بخوارزمية \en{RS256} غير المتماثلة، ومدّة الرمز خمس عشرة دقيقة،
مع رمز تجديد بيتدوّر كل استعمال. وفيه دخول موحّد (\en{Google} و\en{Microsoft})
ومصادقة ثنائية (\en{TOTP}).

\textbf{فصل الأمان —} اتفرّجنا على التدقيق الآليّ: أي أمر متعلّم بـ\code{[Auditable]}
بيتسجّل تلقائيًّا في \code{AuditLog}، وعلى تنقية البيانات الحسّاسة
(\en{PII redaction})، وتحديد المعدّل (\en{rate-limiting})، والتوقيع على أحداث
\en{Stripe}.

\textbf{فصلا الـ\en{AI} —} اتعلّمنا مفاهيم \en{RAG} (استرجاع ثم توليد)
والـ\en{embeddings} والبحث بجيب التمام (\en{cosine similarity})، وبعدها بنينا
المزايا فعليًّا: شات بوت يستند لقاعدة معرفة حيّة، وتوصيات بتحسب درجة حتميّة
(\en{deterministic scoring})، وفحص أهليّة، وكل ده بسقوط آمن ثلاثي المستويات.

\textbf{فصل الاجتماعات والزمن الحقيقي —} شفنا مكالمة الفيديو الكاملة داخل المتصفّح
عبر \en{Azure Communication Services}، من إنشاء الغرفة عند تأكيد الحجز، لإصدار رمز
قصير العمر، للانضمام، للتسجيل على \en{Blob Storage}. وشفنا الـ\en{SignalR} بيدفع
التحديثات لحظيًّا (رسائل، إشعارات، تفاعلات).

\textbf{فصل المدفوعات والمهام —} اتعلّمنا نيّات الدفع (\en{PaymentIntents}) بالتحصيل
اليدوي، وحصّة المنصّة، والمهام الخلفية عبر \en{Hangfire}: تذكير المواعيد، وتسوية
المدفوعات، وكشف الغياب، ومزامنة التحليلات الليلية.

\textbf{هذا الفصل —} لمّينا خطّ التحليلات، وبنلخّص الكتاب، وبنجهّز أسئلة المناقشة.

\begin{summary}[ما الذي يميّز \en{ScholarPath}؟]
لو حد سألك «إيه اللي يخلّي المشروع ده مميّز؟» — دي التلات نقط اللي تقولها بثقة:

\medskip
\textbf{١) تطوير مدفوع بالمواصفات (\en{spec-driven development}).} كل طلب دمج
(\en{PR}) بيشير لمتطلّب وظيفي (\en{FR-xxx}) أو قصّة مستخدم (\en{US-xxx}). كل كيان
له هجرة قاعدة بيانات (\en{migration})، وكل أمر مؤثّر متعلّم بـ\code{[Auditable]}.
المواصفات مش ورق على جنب — هي مصدر الحقيقة اللي الكود بيتبعه.

\medskip
\textbf{٢) ذكاء اصطناعيّ إنتاجيّ بتكلفة حديّة تقترب من الصفر ومع سقوط آمن.}
شات بوت \en{RAG} + توصيات حتميّة + فحص أهليّة، وكلّه وراه \textbf{سقوط ثلاثي
المستويات} بيخلّي النظام يتدهور بلُطف مهما حصل عطل في \en{Azure} — المستخدم دايمًا
بياخد ردّ. وسقف تكلفة يوميّ بيمنع أي فاتورة تطير.

\medskip
\textbf{٣) بنية فيديو متكاملة من الطرف للطرف.} إنشاء غرفة \en{ACS}، وإصدار رمز
قصير العمر، و\en{SDK} في الواجهة، وتسجيل من جهة الخادم على \en{Blob Storage}،
وتحميل \en{MP4} — كلّه مربوط في اثنتي عشرة خطوة مع تتبّع الحضور وكشف الغياب.
\end{summary}

% ═══════════════════════════════════════════════════════════════════════════
\subsection{أسئلة المناقشة المتوقّعة وإجاباتها}
% ═══════════════════════════════════════════════════════════════════════════

الجزء ده هو اللي هيفيدك يوم الدفاع. جمّعنا أكتر الأسئلة المتوقّعة من لجنة المناقشة،
وكل سؤال في صندوقه بإجابة جاهزة ومختصرة. اقراها بصوت عالٍ مرّة أو اتنين قبل الدفاع.

\begin{qa}[ليه \en{Clean Architecture} بدل معمارية طبقية بسيطة؟]
لأنها بتفرض \textbf{عكس الاعتماديّة} (\en{dependency inversion}): طبقة الـ\en{Domain}
ما تعرفش حاجة عن \en{EF Core} ولا \en{Azure}. يعني نقدر نستبدل قاعدة البيانات أو
أي خدمة من غير ما نلمس منطق الأعمال. وأثبتنا ده عمليًّا لمّا حقنّا بدائل
(\en{Stub}) لكل الخدمات الخارجية (\en{ACS}، \en{Stripe}، \en{Azure}) في الاختبارات
وشغّلنا النظام كامل من غيرها.
\end{qa}

\begin{qa}[ليه \en{CQRS} في مشروع تخرّج؟ مش ده أكتر من اللازم؟]
المشروع فيه مئات المتطلّبات الوظيفية موزّعة على خمسة مطوّرين. الـ\en{CQRS} فرض
\textbf{حدود ملكيّة واضحة}: كل مُعالِج ملف \en{C\#} واحد يملكه شخص واحد. من غيره
كان صنف خدمة مشترك هيسبّب تعارضات دمج (\en{merge conflicts}) في كل سبرنت. وكمان
\en{MediatR} إدّانا التدقيق (\en{audit}) كسلوك عرضيّ \textbf{ببلاش}.
\end{qa}

\begin{qa}[ليه ما استخدمتوش قاعدة بيانات متجهات (\en{vector database})؟]
عند حجم يقارب \textbf{692 مستندًا}، المسح الكامل في الذاكرة بجيب التمام
(\en{cosine scan}) بيخلص في أقلّ من عشرة أجزاء من الألف من الثانية. قاعدة المتجهات
هتضيف قفزة شبكة وتكلفة تشغيل من غير أي فايدة في الدقّة عند الحجم ده. ومسار الترقية
موجود: نستبدل \code{KnowledgeRetriever} بتطبيق يعتمد \en{Qdrant} أو \en{pgvector}
لو كبرنا.
\end{qa}

\begin{qa}[هل الـ\en{AI} بيهلوس (\en{hallucinate})؟]
موجّه النظام (\en{system prompt}) \textbf{بيمنع الاختلاق}، والإجابة لازم تستشهد
بكتلة السياق المرقّمة. سياق الـ\en{RAG} مستند لقاعدة البيانات الحيّة؛ ولو مفيش
مستند سجّل درجة فوق العتبة (\en{0.6})، البديل المحلّي بيدّي إجابة عامّة بدل ما
يخمّن. يعني النموذج بيتكلّم من مصدر، مش من خياله.
\end{qa}

\begin{qa}[ليه \en{Azure ACS} للفيديو ومش \en{Jitsi}؟]
\en{Jitsi} بيحتاج خادمًا ذاتيّ الاستضافة بعبء تشغيليّ، ومفيهوش \en{API} تسجيل
مدمج. \en{ACS} بيوفّر خطّ تسجيل مُدار ومتوافق مع خصوصيّة البيانات
(\en{PII-compliant})، وبيتكامل بشكل طبيعيّ مع باقي منظومة \en{Azure} عندنا.
\end{qa}

\begin{qa}[إزاي بتشتغل فوترة الاستشاريين؟]
عبر \en{Stripe PaymentIntents} بتحصيل يدويّ (\en{manual capture}). كارت الطالب
\textbf{يُصرَّح عليه} وقت الحجز، والتحصيل الفعليّ ما بيحصلش إلا لمّا الاستشاري يقبل
الجلسة. وحصّة المنصّة بتتحسب على مستوى المُعالِج عبر \en{Stripe Transfer}.
\end{qa}

\begin{qa}[لو \en{Azure OpenAI} سقط، يحصل إيه؟]
سقوط آمن على \textbf{ثلاثة مستويات}: (١) الإعداد ناقص فبنروح للمحلّي؛ (٢) خطأ
\en{HTTP} في نداء التضمين (\en{embedding}) فبنروح للمحلّي؛ (٣) نداء التوليد فشل
فبنروح للمحلّي. الخدمة المحلّية \code{LocalAiService.AskAsync} بتشغّل نفس استرجاع
الـ\en{RAG} بدون التوليد، وبعدها بترجع لموجّه بالكلمات المفتاحيّة. النتيجة:
\textbf{المستخدم دايمًا بياخد ردّ}.
\end{qa}

\begin{note}
لاحظ إن كل الإجابات دي بتربط بين \textbf{قرار تصميميّ} و\textbf{سبب عمليّ ملموس}.
اللجنة مش بتدوّر على «اخترنا الأحدث»، بتدوّر على «اخترنا الأنسب وليه». خلّي كل
إجابة تبدأ بالسبب.
\end{note}

% ═══════════════════════════════════════════════════════════════════════════
\subsection{الخاتمة}
% ═══════════════════════════════════════════════════════════════════════════

كده نكون قفلنا الدائرة. النظام بيبدأ من طالب بيعمل حساب، بيمرّ باكتشاف المنح
والتقديم والحجز والدفع والاستشارة بالفيديو، وبينتهي ببياناته وهي بتتحوّل لمعرفة في
لوحات التحليل. كل حلقة في السلسلة دي مبنيّة على نفس المبادئ: طبقات نظيفة، أوامر
واستعلامات واضحة، تدقيق آليّ، وسقوط آمن عند أي عطل. ده مش بس نظام شغّال، ده نظام
\textbf{مفهوم} — تقدر تشرح أي جزء منه، وده أهمّ حاجة في يوم الدفاع.

\begin{summary}[الخلاصة في سطور]
\begin{itemize}
  \item خطّ التحليلات بيفصل عالَم \textbf{التشغيل} عن عالَم \textbf{التحليل}:
    \en{SQL Server} ثم \en{ADF} (\en{CDC}, طبقة \en{Bronze}) ثم \en{dbt}
    (\en{Silver} ثم \en{Gold} \en{star schema}) ثم \en{Synapse} ثم \en{Power BI}.
  \item فيه \textbf{فرع لحظيّ} موازٍ: أحداث تُبثّ إلى \en{Event Hubs}، و\en{Stream
    Analytics} يكشف الشذوذ ويُنبّه \en{PagerDuty}.
  \item مبدأ حاكم: التحليلات \textbf{ما تعطّلش التشغيل} — الناشر بيبلع الأخطاء
    ولا يعيد رميها، والبديل الآمن (\en{Stub}) بيشتغل بدون \en{Azure}.
  \item الأحداث المبثوثة بتحمل \textbf{مُعرِّفات وأرقام مجمّعة فقط}، لا بيانات
    شخصيّة حسّاسة.
  \item ما يميّز \en{ScholarPath}: تطوير مدفوع بالمواصفات، \en{AI} إنتاجيّ بتكلفة
    حديّة تقترب من الصفر مع سقوط آمن، وبنية فيديو متكاملة من الطرف للطرف.
  \item في الدفاع: ابدأ كل إجابة بالسبب العمليّ، مش بالتقنية.
\end{itemize}
\end{summary}
